\documentclass[10pt]{article}
\usepackage[margin=1.7cm]{geometry}
\usepackage[T1]{fontenc}
\usepackage{lmodern}
\usepackage{booktabs}
\usepackage{array}
\usepackage{enumitem}
\usepackage{xcolor}
\usepackage[hidelinks]{hyperref}
\setlist{nosep,leftmargin=*}
\newcommand{\pz}{\textbf{[P0]}}
\newcommand{\po}{\textbf{[P1]}}
\newcommand{\pt}{\textbf{[P2]}}
\newcommand{\pth}{\textbf{[P3]}}
\newcommand{\src}[1]{\texttt{#1}}

\title{Independent Review and Revision Plan\\
\large BodhiPromptShield: Pre-Inference Prompt Mediation for Surface-Form Privacy Propagation in LLM Agent Pipelines\\
\normalsize Target venue (authors' decision D1, 6 Sep 2026): PoPETs 2027 Issue 3, deadline 30 Nov 2026 AoE (firm); fallback: ACM TAISAP}
\author{Independent review reset, 5 September 2026}
\date{}

\begin{document}
\sloppy
\maketitle

\section*{Review basis}
This is a fresh review of the manuscript as it stands: the single-file \src{paper/main.tex} (1211 lines; appendix embedded after \verb|\appendix|) and the current \src{paper/main.pdf} (35 pages, \src{acmsmall} journal format, built 3 Sep 2026 from the current source; zero undefined references or citations; three minor overfull boxes). \src{paper/appendix.pdf} is a stale April artifact and was not used. No earlier review notes or completion judgments were consulted for the findings below. After the findings were written, the authors' instruction changed the venue question from ``ICLR 2027'' to ``the best-fitting venue with the highest acceptance probability''; the submission history was then read for venue selection only: the manuscript was rejected by ACM TOPS (submission TOPS-2026-06-0395) on an Associate Editor recommendation with six comments, after an earlier CCS cycle. Those six comments converge on the same defects identified independently here; the AE's closing sentence---``a clear technical contribution and a new evaluation on real agent stacks with strong baselines and realistic privacy attacks''---is the bar every venue below will apply.

\section*{Venue recommendation}
\begin{center}
\small
\begin{tabular}{@{}p{0.19\textwidth}p{0.23\textwidth}p{0.26\textwidth}p{0.24\textwidth}@{}}
\toprule
Venue & Fit & Format / deadline (verified 5 Sep 2026) & Assessment \\
\midrule
\textbf{ACM TAISAP} (Trans.\ on AI Security and Privacy) & Scope is exactly ``privacy concerns related to AI model \ldots deployment'' and AI-enabled systems; theoretical and applied work welcome & ACM journal, \src{acmart} (already used); no deadline; open access with no charge 2026--2029; new journal opened for submissions in 2025. Confirm the length cap and blinding policy on the author-guidelines page (the ACM norm, as at TOPS, is 35 pages total) & \textbf{Fallback} (authors chose PoPETs on 6 Sep 2026). Least reformatting, keeps the systems/threat-model framing, allows the full appendix, and a young journal has a shorter queue. The AE's six points must still be answered with new evidence. \\
\addlinespace
\textbf{PoPETs 2027 Issue 3/4} (PETS) & Privacy engineering and ``machine learning and privacy'' are listed topics; the community values threat models, honest scope, and artifacts & 12 typeset main-body pages (13 on revise), excluding ethics, open-science, AI-use sections, acknowledgements, bibliography, appendices; double-blind; decisions are accept / revise / reject with a four-month editor-guided revision path; rejected papers skip one issue. Issue 3 deadline 30 Nov 2026, Issue 4 deadline 28 Feb 2027 (firm) & \textbf{Selected primary (D1).} Best conference fit; the ``revise'' outcome is a realistic second chance. Requires a $\approx$55\% cut of the body and full anonymization. \\
\addlinespace
USENIX Security '27 Cycle 2 & Applied security; agent privacy is topical & Registration 19 Jan 2027, submission 26 Jan 2027 & Stretch target only if E1--E7 (below) yield a clear head-to-head win. \\
\addlinespace
ICLR 2027 & Weak: no learning contribution; reviewers would discard the systems material & 9 pages, 25 Sep 2026 & Not recommended. \\
\addlinespace
ACM TOPS resubmission & In scope, but just rejected & 35 pages & Not recommended without EiC invitation. \\
\bottomrule
\end{tabular}
\end{center}
The experimental work is identical for all rows. Decision taken 6 Sep 2026: PoPETs Issue 3 is the target, so the 30 November deadline and the 12-page main body are hard constraints and every ``PoPETs delta'' below is mandatory; the TAISAP notes are retained only as the fallback if E1--E7 slip past Issue 4 (28 Feb 2027). CPPB is demoted to the appendix (D2); no rebuild this cycle.

\section*{Summary of the submission}
The paper argues that prompt privacy in LLM agent pipelines should be treated as a propagation-control problem across retrieval, memory, and tool boundaries rather than as one-shot document redaction. It proposes a mediation layer (span detection, policy routing to typed placeholders / semantic abstraction / secure symbolic mapping, late authorized restoration), defines PER and stage-wise SPE, and evaluates on a paper-specific 256-prompt controlled benchmark (CPPB), on AI4Privacy text pushed through a reference agent, on 493 PrivacyLens trajectories, on TAB, and on OCR slices. The best-supported findings are: (i) with a learned detector, exact-identifier survival across boundaries falls to PER 7.4\% (AI4Privacy) and 1.8\% (PrivacyLens); (ii) 83\% of PrivacyLens sensitive facts remain semantically conveyed after mediation; (iii) the released heuristic detector is trivially evaded (zero-width ASR 99.7\%) while the learned detector resists typographic but not paraphrase attacks (28.2\%); (iv) a per-session mapping design confines compromise to one session except under master-seed loss.

\section*{Overall assessment}
\textbf{Format is not the blocker for TAISAP; the merits are.} The manuscript would most likely be rejected again, for four reasons reviewers will state almost verbatim (and the TOPS AE already did):

\begin{enumerate}
\item \textbf{The headline comparison is a loss the paper concedes.} Enterprise staged redaction has lower direct exposure (PER 8.1\% vs 9.3\%), the utility gain is $+0.02/+0.02$ on an in-house rubric, and the appendix blind review (Table~\src{tab:blindutility}) reports task success of 12--13.5\% with 2--9 of 148 outputs judged useful. A reviewer will read 0.92 TSR in the main text and 12\% TSR in the appendix and stop trusting the utility numbers.
\item \textbf{The controlled benchmark is not reproducible from the release and its central table is not measured.} The main text states the released CPPB manifest contains ``lightweight prompt stubs'' and the score logs are ``historical/generated records rather than a fresh end-to-end model rerun''; the appendix caption of Table~\src{tab:propagation} says ``Numbers are a controlled benchmark surface, not measured system output.'' Unreproducible headline tables are fatal at any of the venues above.
\item \textbf{The evidence bar is below 2026 norms.} Downstream models are \src{llama3:latest} and a 3.2B Q4 \src{llama3.2}; the learned detector is a 17-feature logistic regression; there is no LLM-based sanitization baseline, no real (installed, version-pinned) Presidio, no Casper or PrivacyRestore re-implementation, no confidence intervals on any headline table, and no task-native utility metric on the real agent benchmarks.
\item \textbf{The manuscript argues against itself.} Roughly half of the results prose is scope disclaimers, the Limitations section enumerates more than a dozen open gaps, and captions announce that their own numbers are not evidence. Combined with pervasive revision-history language (``the latest review required'', ``the new audit'', ``moved to the companion appendix'', a section titled ``Revision-Cycle Live Pilot'', an artifact filename in the appendix), the paper reads as an internal audit report rather than a research paper.
\end{enumerate}

\textbf{What would make it acceptable at a security/privacy venue.} The systems framing is an asset here and should stay. What must change is the contribution and its evidence: (a) the mechanism must be compared head-to-head, under one protocol, against Casper-style sanitization, real Presidio, and an LLM sanitizer, and shown to control something they do not---restoration timing and cross-boundary propagation on \emph{real} agent stacks with strong models; (b) the propagation-aware protocol (PER, SPE, and a semantic-leakage measure) must be applied to real trajectories with uncertainty; (c) the 83\% semantic-conveyance result must be turned from a caveat into a measured finding across all methods, with the abstraction mode evaluated as a mitigation; (d) the writing must state one bounded claim confidently and its limitations once. The AE's phrase---``a new evaluation on real agent stacks with strong baselines and realistic privacy attacks''---is the specification; A2 below implements it.

\section{Findings}
Priorities: \pz\ blocks submission; \po\ accept-critical; \pt\ materially strengthens; \pth\ camera-ready or optional.

\subsection{Compliance and formatting}
\begin{enumerate}[label=F\arabic*.]
\item \pz\ \textbf{Journal metadata and length.} \src{main.tex} lines 16--24 carry TOPS metadata (\verb|\acmJournal{TOPS}|, volume/number/article placeholders). For TAISAP, swap the journal code and confirm the length cap on the author-guidelines page; at exactly 35 pages the paper has no headroom for the new results, so plan the rewrite for $\le 30$ pages (structure plan below). \emph{PoPETs (measured 6 Sep 2026):} the manuscript has been converted to the PoPETs 2027 template (\src{acmart} \src{sigconf,anonymous,review} + \src{popets.sty}); it initially built at 21 pages with the main body ending on page 16; after the F12 cleanup and the restructure of 6 Sep 2026 (CPPB material to Appendix~B, disclaimed tables and the blind-review section removed, a real-data Datasets subsection) it builds at 20 pages with the main body ending on page 12, i.e.\ within the limit before the new results are added.
\item \pz\ \textbf{Anonymity.} Author names, affiliations, and e-mails at lines 57--75; \verb|\shortauthors{Ma et al.}|; an \src{acks} block; and the URL \url{https://github.com/mabo1215/BodhiPromptShield} in the abstract (line 78) and the reproducibility paragraph (line 527). Confirm TAISAP's blinding policy; if double-blind, all of these must go and the repository must be mirrored anonymously (\url{anonymous.4open.science} or an anonymized zip). \emph{PoPETs delta:} double-blind is mandatory; any identity leak is a desk rejection.
\item \pz\ \textbf{Required statements.} PoPETs 2027 mandates three named sections from its template---\emph{ethical considerations}, \emph{open science}, and \emph{AI use}---none counted toward the 12 pages (author instructions, fetched 6 Sep 2026). The ethics section must state that PhysioNet-derived clinical text was used under credentialing and that TAB is public ECHR case text, and that no real user data was collected. The same instructions require that \textbf{the first page clearly state how the work is relevant to real-world privacy applications}; the current abstract and introduction do not do this explicitly. Template: \src{submission-template-2027.zip} from the PoPETs authors page; fonts must be Type1/TrueType, embedded and subset; content confined to the two-column margins.
\item \pz\ \textbf{Legacy cross-reference.} Line 559: ``as described in Section~IV'' is a hard-coded Roman numeral left from an earlier template; the framework section is Section~4 (\verb|\ref{sec:framework}|).
\item \po\ \textbf{Bibliography style.} \verb|\bibliographystyle{ACM-Reference-Format}| (line 915) is correct for TAISAP. \emph{PoPETs delta:} switch to the template's style and citation macros.
\item \po\ \textbf{Promised figures that do not exist.} Line 511 (``A companion appendix figure further visualizes \ldots''), line 701 (``supporting restoration-boundary table and visualization''), and line 733 (``supporting ablation table and trade-off plot'') refer to figures that are not in the PDF. The PDF has exactly two figures (\src{fig:workflow}, \src{fig:operating\_points}); four unused PNGs sit in \src{paper/figs/} (\src{agent\_pipeline\_summary}, \src{agent\_propagation\_curves}, \src{cppb\_benchmark\_composition}, \src{restoration\_ablation\_tradeoffs}). Either include them or delete the sentences.
\item \po\ \textbf{Dangling labels and uncited references.} Labeled but never referenced: \src{fig:workflow} (the architecture figure is never cited in text), \src{tab:cppb\_audit}, \src{tab:mapping\_surface}, \src{tab:pi\_profiles}, \src{sec:blind\_utility}. Bib entries never cited: \src{shao2024privacylens} (PrivacyLens is a headline data source and its paper is never cited---an attribution failure reviewers will flag), \src{jaume2019funsd} (FUNSD is used in the OCR slices), \src{debenedetti2024agentdojo}.
\item \pt\ \textbf{Stale build layout.} \src{paper/appendix.pdf}, \src{paper/build/appendix.*}, and \src{paper/AGENTS.md} still describe a separate appendix file; the manuscript is single-file. Stop shipping the stale appendix PDF.
\end{enumerate}

\subsection{Contribution, positioning, and narrative}
\begin{enumerate}[label=F\arabic*.,resume]
\item \po\ \textbf{Contribution list is not defensible as written.} Of the four bullets in the Introduction, (2) the mediation mechanism is conceded to be prior art in \S3.3 (``This framing does not make the underlying replacement operation new''), and (4) ``making the evidence scope explicit'' is not a contribution. The AE said explicitly that ``the propagation framing alone does not constitute sufficient novelty.'' The defensible contribution set for a security venue is: (1) the propagation-control formulation with a threat-to-evaluation mapping; (2) restoration governance (per-session mapping + boundary-timed restoration) as a \emph{measured} mechanism, compared head-to-head with Casper-style, Presidio, and LLM-sanitizer baselines on real agent stacks; (3) the propagation-aware protocol with PER, SPE, and a semantic-leakage metric, and the findings it produces (identifier propagation is nearly eliminated, semantic conveyance is not, late restoration cuts boundary leakage by a measured factor).
\item \po\ \textbf{The abstract leads with a loss and no finding.} The current abstract gives eight numbers, two of which are unfavorable, and ends with three sentences of disclaimers. Rewrite to state the threat, the mechanism, the protocol, the three findings, and one sentence of scope.
\item \po\ \textbf{Self-undermining prose.} Recurrent constructions such as ``we make no \ldots claim'', ``not evidence of'', ``cannot establish'', ``should not be interpreted as'' appear in nearly every results paragraph. State each limitation once, in the Limitations section, and let results paragraphs report what was measured.
\item \po\ \textbf{Revision-history and repository language must go.} Examples: ``The span-level replacement analysis from an earlier version has been removed'' (\S2.4), ``the latest review required a stricter provenance distinction'' (\S6.8), ``the new metadata audit flags'' (\S6.10), ``The appendix now reports'' (\S6.5, \S6.7), ``is moved to the companion appendix'' (\S6.3), ``The repository now distinguishes'' (\S5.7), appendix section ``Revision-Cycle Live Pilot and Semantic Probe'' and Table~\src{tab:revision\_execution} (``Revision-cycle execution audit''), and the artifact filename \src{review\_revision\_audit.csv} in Appendix~A. None of this is legible to a first-time reader.
\item \po\ \textbf{Tables whose captions disclaim their own content.} Table~\src{tab:propagation} (``not measured system output''), Table~\src{tab:crossmodel} (``not evidence of cross-model portability''), Table~\src{tab:cppb\_audit} (``not promoted to live held-out evidence''). Delete these tables or replace them with measured results.
\item \pt\ \textbf{Related work under-covers the agent-privacy literature.} Missing or unengaged: PrivacyLens itself (Shao et al.\ 2024, in the bib but uncited), ConfAIde (Mireshghallah et al.\ 2024, contextual-integrity leakage), AgentDojo (in bib, uncited), PAPILLON / privacy-conscious delegation (Siyan et al.\ 2024), Rescriber (Zhou et al.\ 2025), LLM-based anonymizers and their adversaries (Staab et al.\ 2024), Dou et al.\ 2024 on reducing self-disclosure risk, and transformer PII detectors trained on AI4Privacy. Add 6--10 citations with one sentence each; this also supplies the missing baselines.
\end{enumerate}

\subsection{Experimental rigor}
\begin{enumerate}[label=F\arabic*.,resume]
\item \po\ \textbf{Downstream models are too weak to support utility claims.} \S5.7: ``local Ollama \src{llama3:latest}'' and ``\src{llama3.2:latest} \ldots 3.2B, Q4\_K\_M''. Utility must be measured with at least one strong open model (Llama-3.1-70B-Instruct or Qwen2.5-72B/Qwen3, named with revision) and one API model (named with version string), each at fixed decoding. This also answers AE comment 5 (withheld model identities).
\item \po\ \textbf{Missing baselines (AE comment 1).} (a) An LLM zero-shot/few-shot sanitizer (``rewrite this prompt removing PII, keep task intent'')---the most obvious 2026 baseline and absent; (b) real Microsoft Presidio, installed and version-pinned, instead of the ``Presidio-class approximation''; (c) a Casper-style re-implementation (rules + NER + dummy values + local restore; a few hundred lines) run under the same protocol; (d) an off-the-shelf transformer PII detector (e.g., a GLiNER-PII checkpoint or the public AI4Privacy-trained token classifier) as the learned detector, replacing or accompanying the 17-feature logistic regression (precision 0.36), which reviewers will call a straw learned baseline. PrivacyRestore (activation steering, server-side) can be discussed rather than reproduced if the architectural reason is stated. In addition, the row labelled ``NER-only masking'' in every table is not a named-entity recognizer: the released code implements it as capitalization, title, and location-context regular expressions. Either relabel it (``capitalization heuristic'') or replace it with a real NER baseline (spaCy \src{en\_core\_web\_lg} entities); reviewers who open the artifact will treat the current label as misreporting.
\item \po\ \textbf{No uncertainty on any headline number.} Tables~\src{tab:measprop}, \src{tab:per}, \src{tab:utility}, \src{tab:advmeas}, \src{tab:pi\_sensitivity} report point estimates only. Add bootstrap 95\% CIs over documents/trajectories for all exposure metrics and $\geq 3$ seeds or paired bootstrap for utility; report paired differences vs.\ the strongest baseline.
\item \po\ \textbf{SPE monotonicity claim is contradicted by the paper's own data.} \S4.5 states $\mathrm{SPE}_{k+1}\le\mathrm{SPE}_k$ along retrieval$\to$memory$\to$tool. Table~\src{tab:measprop}: No-protection on AI4Privacy is 48.7 / 100.0 / 74.8; on PrivacyLens every method has Tool $=0.0$ while Retrieval is 0 or 100. Either the surfaces are not sequential in the reference agent (then define SPE per surface, not as a decay), or the Tool column for PrivacyLens has no gold spans (then say so and do not report zeros as suppression). Also explain why NER-only masking is 100/98.9 on PrivacyLens.
\item \po\ \textbf{Utility metric is not credible (AE comment 4).} AC/TSR are an unvalidated in-house rubric; the appendix blind review gives 12--13.5\% task success and 2--9/148 useful outputs, and per-template TSR ranges from 0\% to 50\%. Replace with task-native metrics on the real benchmarks (PrivacyLens helpfulness judgment; AgentDojo utility / task success with tool calls), an LLM-judge rubric with a reported human-agreement subset ($n\ge 100$, Cohen's $\kappa$), and report both under each mediation method. Retire AC/TSR from the main text.
\item \po\ \textbf{CPPB cannot carry headline tables (AE comment 2).} Four of the seven main-text tables (\src{tab:per}, \src{tab:utility}, \src{tab:pi\_sensitivity}, \src{tab:latency}) rest on CPPB logs that the paper says cannot be regenerated; the released reconstruction uses a single injected value per category (one person name for all 32 person-name prompts), so a memorized string would score perfectly. Resolutions: (i) rebuild CPPB with diverse generated values (multiple surface forms per category, realistic distractors), run it live with named models, and keep it as a controlled ablation surface; or (ii) demote CPPB to the appendix and promote the real-data results. Resolved 6 Sep 2026 (D2): (ii), demote to the appendix.
\item \po\ \textbf{Restoration-timing evidence is controlled, not measured.} The paper's most distinctive claim (late restoration: BLR 1.8\% vs 9.7\%, Table~\src{tab:restore}) is on the controlled surface. Measure it on real tool-calling tasks where an exact value is required by a tool (AgentDojo has such tasks; PrivacyLens trajectories include tool actions): report TSR, RSR, and BLR for no / early / late restoration with CIs. This is the mechanism-level contribution the AE asked for.
\item \po\ \textbf{Semantic leakage should be a result, not a limitation (AE comment 3).} The 83\% ``semantically conveyed'' finding is the most interesting number in the paper and is introduced as a caveat. Measure it for every method (raw, regex, NER, Presidio, LLM-sanitizer, Casper-style, placeholder, abstraction, symbolic) on all 493 PrivacyLens trajectories with an LLM-judge ``is fact $f$ still inferable from the released surface?'' protocol plus a human-checked subset; report the surface-form vs.\ semantic leakage gap as a headline figure, and evaluate the abstraction mode as a mitigation with paired utility. The 12-probe mitigation panel (\S6.6, \S6.13) is too small for the main text.
\item \pt\ \textbf{Internal numeric inconsistencies.} (a) Regex-only / NER-only PER is 61.3 / 48.7 in Table~\src{tab:per} but 61.6 / 49.0 in Table~\src{tab:judge} for the same CPPB baselines. (b) Learned detector in-domain: F1 0.53, recall 0.95, PER 7.6\% (Table~\src{tab:lightclf}) vs.\ F1 0.51, recall 0.96, PER 7.4\% (Table~\src{tab:crossdom}). (c) \S5.7 gives the latency hardware (Ryzen 7 3700X, 32\,GB, Windows 11) while \S6.12 says ``the current snapshot does not yet bundle the hardware manifest''. (d) Homoglyph: ASR 65.4\% in Table~\src{tab:advmeas} vs.\ ``lowers average homoglyph exposure from 94.1\% to 43.9\%'' in \S6.14, with no reconciliation. (e) The eighth CPPB category is ``Context-dependent confidential spans'' in Table~\src{tab:cppb\_card} but ``image-derived text'' in Appendix~B. (f) Context-inference accuracy is reported three times with different panels (100$\to$50; 50/25/8.3/16.7). Reconcile every pair or remove one member.
\item \pt\ \textbf{Adversarial evaluation is one-sided.} Table~\src{tab:advmeas} reports evasion of the detector only; also report the downstream effect (does an evaded span then propagate? SPE under attack) and include an adaptive attacker with knowledge of the routing policy (restoration-trigger injection success), which the threat model lists but the results do not measure. This is the ``realistic privacy attacks'' item in the AE's closing sentence.
\item \pt\ \textbf{Formal-scope paragraph (AE comment 6).} \S2.4 still discusses the removed LDP analysis. Reduce to two sentences stating that the system is deterministic and offers no formal guarantee; do not narrate what was removed.
\item \pt\ \textbf{OCR slice.} The 64-prompt CPPB OCR subset and the CORD/FUNSD/SROIE runs are asserted (\S6.9) with three numbers and no table. Add a compact appendix table with the OCR engine and version, or remove the paragraph. FUNSD must be cited if used.
\item \pth\ \textbf{Cross-model portability.} Table~\src{tab:crossmodel} (alias-only backends) should be deleted. If a portability claim is wanted, run the final protocol on the two named models from F15 and report it directly.
\end{enumerate}

\subsection{Traceability to the TOPS Associate Editor comments}
\begin{center}
\small
\begin{tabular}{@{}p{0.06\textwidth}p{0.42\textwidth}p{0.46\textwidth}@{}}
\toprule
AE & Comment (abridged) & Answered by \\
\midrule
1 & Casper / PrivacyRestore not cited or compared; propagation framing alone is not novelty & F9, F16(c), E1 (Casper-style + Presidio + LLM sanitizer under one protocol), E4 (restoration governance as measured mechanism) \\
2 & CPPB: 256 prompts / 32 templates; full-dataset results; identical dev/test & F20, E8 (rebuild with diverse values and live rerun, or demote) \\
3 & Deterministic reference agent; 83\% semantic exposure; 50\% inference accuracy after sanitization & F18, F22, E5 (semantic-leak gap across methods, abstraction as mitigation), E7 \\
4 & Enterprise redaction beats proposed on PER; +0.02 utility on an unvalidated rubric & F19, E3 (task-native utility with strong models and CIs), abstract/contribution rewrite (F9--F10) \\
5 & Detector F1 0.53 / precision 0.36; model identities withheld; OCR/runtime records missing & F15, F16(d), F26, E2, named models throughout; run manifests in the artifact \\
6 & LDP analysis irrelevant to the evaluated system & F25 \\
\bottomrule
\end{tabular}
\end{center}

\section{Target paper structure}
\subsection*{TAISAP (ACM journal, target $\le 30$ pages \src{acmsmall})}
Keep the current section order. Cuts: merge \S2.1--2.3 and drop the design-objective equation and \S2.4 to two sentences ($-2$ pages); Related Work to 1.5 pages ($-1$); Framework: keep overview figure, three modes, one routing paragraph with Table~\src{tab:pi\_rules}, mapping design with Table~\src{tab:mapcompromise}; move \src{tab:pi\_profiles}, \src{tab:mapping\_surface}, the $\alpha/\beta$ equations, and the algorithm to the appendix ($-3$); Setup: rebuild around real benchmarks, named models, and the E1--E2 baseline roster ($-2$); Results: five measured results (E3, E4, E5, propagation with CIs, adversary with CIs), two new figures (propagation curves on real agents; surface-vs-semantic leakage per method) ($\pm 0$ net after removing disclaimers); delete \src{tab:cppb\_audit}, \src{tab:propagation}, \src{tab:crossmodel}, \src{tab:revision\_execution}, \src{tab:blindutility}, \src{tab:blindutilitytemplate} ($-3$); Limitations to one paragraph. Net: $\approx 35 \to 27$ pages before the new results, $\approx 30$ after.

\subsection*{PoPETs delta (12 typeset main-body pages)}
\begin{center}
\begin{tabular}{@{}p{0.6\textwidth}r@{}}
\toprule
Section & Pages \\
\midrule
Abstract; 1 Introduction with three contributions & 1.3 \\
2 Problem setting and threat model (4-row threat table) & 1.2 \\
3 Related work & 0.8 \\
4 Mediation mechanism (figure; modes; routing paragraph; restoration boundary; mapping design in 6 lines) & 2.0 \\
5 Evaluation protocol (metrics incl.\ semantic leak; datasets; baselines; models) & 1.7 \\
6 Results (propagation on real agents; utility with strong models; semantic-leak gap; restoration timing; adversary) & 4.0 \\
7 Limitations and conclusion & 0.7 \\
Ethics, Open science, AI use, references, appendix & uncounted \\
\bottomrule
\end{tabular}
\end{center}

\section{Revision actions}
Owners: \emph{paper} = manuscript edits in \src{paper/}; \emph{src} = code and experiments in \src{src/} with figures to \src{paper/figs/}. Days are calendar days from 5 September. GPU work must start on day 1 and overlap with the paper restructuring; every long-running script writes incrementally and resumes.

\subsection*{A1. Compliance and hygiene (paper; days 1--4) \pz}
\begin{enumerate}[label=A1.\arabic*]
\item Swap the journal metadata from TOPS to TAISAP; grep for ``TOPS'' and remove every occurrence. Confirm the length cap and blinding policy on the TAISAP author-guidelines page and record both in \src{docs/progress.md}.
\item If double-blind: remove the author block, \verb|\shortauthors|, and \src{acks}; replace both GitHub URLs with an anonymous mirror; grep for ``Auckland'', ``AUT'', ``Guilin'', author surnames, and e-mail domains.
\item Add Ethics, Open-science/artifact, and AI-use statements.
\item Fix ``Section~IV'' (line 559); reference \src{fig:workflow} in \S4.1; cite \src{shao2024privacylens} where PrivacyLens is first used, \src{jaume2019funsd} in the OCR paragraph, and \src{debenedetti2024agentdojo} once AgentDojo is used.
\item Delete the three sentences promising nonexistent appendix figures (lines 511, 701, 733) or include the corresponding figures.
\item Remove every revision-history and repository phrase listed in F12, and the filename in Appendix~A; delete the appendix pilot section per the structure plan.
\item Reconcile all six inconsistencies in F23.
\item \pz\ PoPETs requires that resubmissions of papers rejected \emph{from any venue with reviews} include those prior reviews in the submission. Attach the TOPS Associate Editor comments (\src{docs/AESuggestion.txt}, anonymized) together with a point-by-point response built from the traceability table above.
\end{enumerate}

\subsection*{A2. Experiments (src; days 1--45) \po}
\begin{center}
\small
\begin{tabular}{@{}lp{0.30\textwidth}p{0.22\textwidth}p{0.20\textwidth}r@{}}
\toprule
ID & Purpose & Data / models & Deliverable & Days \\
\midrule
E1 & LLM zero-shot sanitizer baseline; real Presidio (pinned); Casper-style baseline & AI4Privacy test, PrivacyLens 493, TAB; Qwen2.5-7B-Instruct or Llama-3.1-8B for the sanitizer & new rows in propagation and semantic-leak tables & 4 \\
E2 & Off-the-shelf transformer PII detector as the learned detector & GLiNER-PII or public AI4Privacy token classifier & replaces LR rows; detector table in appendix & 2 \\
E3 & Utility with strong downstream models under each method & PrivacyLens helpfulness; AgentDojo utility/task success; one $\ge$70B open model + one API model, named & main utility table with CIs & 7 \\
E4 & Restoration timing on real tool-calling tasks & AgentDojo exact-value tool tasks; PrivacyLens tool actions & main table: TSR/RSR/BLR for none/early/late with CIs & 4 \\
E5 & Semantic-leak gap across all methods; abstraction as mitigation with paired utility & all 493 PrivacyLens trajectories; LLM-judge + $n{=}100$ human-checked subset with $\kappa$ & headline figure: surface vs.\ semantic leakage per method & 6 \\
E6 & Bootstrap CIs and paired tests for every retained table; 3 seeds where stochastic & all of the above & CI columns; significance notes & 2 \\
E7 & Adaptive attacker with downstream effect (SPE under attack; restoration-trigger injection) & AI4Privacy 200-doc slice; PrivacyLens subset & extended adversary table & 3 \\
E8 & Demote CPPB to the appendix (D2): move \src{tab:per}, \src{tab:utility}, \src{tab:pi\_sensitivity}, \src{tab:latency} and the CPPB card there; delete the disclaimed CPPB tables & CPPB & appendix section & 1 \\
\bottomrule
\end{tabular}
\end{center}
Record model digests, decoding settings, seeds, and hardware in run manifests; keep all identifiers out of the manuscript prose.

\subsection*{A3. Rewrite (paper; days 10--55) \po}
\begin{enumerate}[label=A3.\arabic*]
\item New abstract and introduction per F9--F10 with the three contributions listed in F9.
\item Merge \S2 and the threat model; 4-row threat table; two-sentence formal-scope statement.
\item Compress \S4 to the mechanism description; move routing profiles, mapping-surface table, coefficient equations, and the algorithm to the appendix; keep the compromise curve.
\item Rebuild \S5 around the real benchmarks, named models, and the E1--E2 baseline roster; define every metric once, including the semantic-leakage metric.
\item Rebuild \S6 around five measured results (E3, E4, E5, propagation with CIs, adversary with CIs); each paragraph states setup, number, and conclusion---no disclaimers.
\item One Limitations paragraph (semantic inference remains open; no formal guarantee; benchmarks are synthetic-user or reconstructed; single-language; OCR bounded). One conclusion paragraph.
\item Appendix: CPPB card, algorithm, routing tables, OCR table with engine version, TAB table, policy sweep, per-category tables, full experimental details.
\item Build; verify the length cap, zero undefined references, no identity strings if double-blind, every figure and table referenced.
\end{enumerate}

\subsection*{A4. Submission logistics \pz}
\begin{enumerate}[label=A4.\arabic*]
\item Deadline is fixed: PoPETs Issue 3, 30 November 2026 AoE (firm). Checkpoint 31 October: if E1--E7 are not complete, fall back to Issue 4 (28 February 2027) rather than submit thin evidence; a rejected PoPETs paper must skip one issue.
\item Prepare the artifact (anonymized if required) with pinned environments and the run manifests from A2.
\item Cover letter: state that the paper was previously reviewed at TOPS and list, per AE comment, the new evidence (use the traceability table above).
\end{enumerate}

\section{Decisions required from the authors}
\begin{enumerate}[label=D\arabic*.]
\item \textbf{Venue.} \emph{Resolved 6 Sep 2026:} PoPETs 2027 Issue 3 (30 November deadline, 12-page main body); TAISAP is the fallback.
\item \textbf{CPPB disposition.} \emph{Resolved 6 Sep 2026:} demote to the appendix; no rebuild this cycle.
\item \textbf{Model budget.} \emph{Host inventory, 6 Sep 2026:} the local workstation has an RTX 3070 (8\,GB) and only the cloud-proxied \src{kimi-k2.5:cloud} in Ollama; vGPU 3090 is reachable but in no-GPU mode, its Ollama store holds only \src{llama3.2:latest} (1.9\,GB), and its shared storage holds vision models only (Qwen3-VL-8B/4B, EGM-4B, LocateAnything-3B, SAM3); H800, PRO 6000, and A800 are powered off; the 4c 3090 is not reachable from this network. \textbf{No $\ge$70B text model exists on any host today.} \emph{Recommended plan (sizes verified on the Ollama library):} power on the H800 or A800 (80\,GB) and pull \src{llama3.3:70b} (43\,GB, Q4\_K\_M) as the strong open downstream model, plus \src{qwen3:32b} (20\,GB) as the second open model for the portability row; both fit one 80\,GB card. Use an already-licensed API model (the OpenAI or Claude account used for the blind ratings, with its exact version string) as the API downstream model, with \src{kimi-k2.5:cloud} (digest-pinned) as fallback. For E5/E3 judging, use a different family from the downstream model (e.g., \src{qwen3:32b} on the 3090 once its GPU is enabled) plus the $n{=}100$ human subset. \src{qwen3:8b} (5.2\,GB) serves as the E1 LLM-sanitizer baseline. Check \src{df -h} and \src{df -i} on the target host before pulling 43\,GB. \emph{Still open:} which host to power on, and the API spend ceiling for judging.
\end{enumerate}

\end{document}
